\chapter{RESULTS AND EVALUATION}

This chapter presents the experimental results. It first describes the evaluation setup and metrics, then compares all models on a common test set, quantifies the differences with statistical tests, analyzes performance under different operating conditions, demonstrates generalization across routes, compares the system with the reference study, and concludes with the methodological findings.

\section{Experimental Setup}

The main experiments are conducted on route 502 (Cengizhan -- Halkapınar Metro), the most intensively collected line. After cleaning, the dataset for this route contains $81{,}575$ stop-to-stop segments gathered from $48$ distinct buses over a window spanning April to June 2026, during service hours of roughly 06:00--22:00. The mean segment travel time is $1.21$~minutes with a standard deviation of $1.54$~minutes, confirming that inter-stop travel times are short and highly variable (a coefficient of variation above one). The data are split chronologically into $65{,}260$ training and $16{,}315$ test segments.

Four standard regression metrics are used. For $n$ test segments with true travel times $y_i$ and predictions $\hat{y}_i$:
\begin{align}
\text{MAE} &= \frac{1}{n}\sum_{i=1}^{n}\lvert y_i - \hat{y}_i\rvert, &
\text{RMSE} &= \sqrt{\frac{1}{n}\sum_{i=1}^{n}(y_i - \hat{y}_i)^2}, \\
\text{MAPE} &= \frac{100}{n}\sum_{i=1}^{n}\left\lvert\frac{y_i - \hat{y}_i}{y_i}\right\rvert, &
R^2 &= 1 - \frac{\sum_i (y_i-\hat{y}_i)^2}{\sum_i (y_i-\bar{y})^2}.
\end{align}
MAE is the primary metric, as it is robust and directly interpretable in minutes; $R^2$ measures how much of the variance the model explains.

\section{Overall Model Comparison}

Table~\ref{tab:overall} reports the performance of all models on the same route-502 test set. The two non-learned baselines are clearly outperformed by the learned models: the naïve schedule-based predictor and the historical average reach MAE values of $0.73$ and $0.65$~minutes respectively, whereas the improved learned models cluster around $0.43$~minutes. Among the learned models, XGBoost achieves the lowest MAE ($0.4327$), the dual-input LSTM achieves the best $R^2$ and RMSE, and the Random Forest is marginally behind both.

\begin{table}[!htbp]
    \centering
    \caption{\label{tab:overall}Model performance on the route-502 test set (same test set for all models). Lower is better for MAE/RMSE/MAPE; higher is better for $R^2$.}
    \begin{tabular}{@{}lrrrr@{}}
        \toprule
        \textbf{Model} & \textbf{MAE (min)} & \textbf{RMSE (min)} & \textbf{MAPE (\%)} & \textbf{$R^2$} \\
        \midrule
        \textbf{XGBoost (Improved)} & \textbf{0.4327} & 0.903 & 41.1 & 0.626 \\
        LSTM (dual-input)          & 0.4345 & \textbf{0.891} & 40.7 & \textbf{0.636} \\
        Random Forest (Improved)   & 0.4378 & 0.896 & 41.6 & 0.633 \\
        Historical Average         & 0.645  & 1.317 & 62.7 & 0.206 \\
        Naïve (GTFS schedule)      & 0.734  & 1.516 & 68.2 & $-0.05$ \\
        \bottomrule
    \end{tabular}
\end{table}

A subtle but important point concerns the LSTM. Evaluated in isolation on its own sequence-level test set (which, by construction, excludes the harder cold-start segments), the LSTM reaches an MAE of about $0.35$~minutes. This figure is \emph{not} directly comparable with the tree models, because it is measured on an easier subset. When the LSTM is evaluated on exactly the same full segment-level test set as the other models, its MAE is $0.4345$~minutes, as shown in Table~\ref{tab:overall}. Reporting the same-test-set figure is essential for a fair comparison and is the basis for the statistical analysis below.

\section{Statistical Significance Testing}

To determine whether the differences between the learned models are real or within noise, paired statistical tests (paired $t$-test and Wilcoxon signed-rank test) were applied to the per-segment absolute errors on the common test set. The results are summarized in Table~\ref{tab:sig}.

\begin{table}[!htbp]
    \centering
    \caption{\label{tab:sig}Paired significance tests on the common test set ($p<0.05$ indicates a statistically significant difference).}
    \begin{tabular}{@{}lrl@{}}
        \toprule
        \textbf{Comparison} & \textbf{$p$ ($t$-test)} & \textbf{Significant?} \\
        \midrule
        XGBoost vs. LSTM            & 0.38   & No \\
        XGBoost vs. Random Forest   & 0.0055 & Yes \\
        LSTM vs. Random Forest      & 0.0014 & Yes \\
        XGBoost vs. Historical Avg. & $<10^{-3}$ & Yes \\
        XGBoost vs. Naïve (GTFS)    & $<10^{-3}$ & Yes \\
        \bottomrule
    \end{tabular}
\end{table}

The tests yield a clear ordering, $\text{XGBoost} \approx \text{LSTM} > \text{RF}$. XGBoost and the LSTM are statistically \emph{indistinguishable} ($p=0.38$): the deep model does not outperform the classical one. Both, however, are significantly better than the Random Forest, and all learned models are overwhelmingly better than the baselines. This result corrects an earlier, informal conclusion that the LSTM was the single best model; that impression arose only because the LSTM had been evaluated on an easier test subset. On a fair, common test set, the simpler XGBoost is the most practical choice, since it matches the deep model in accuracy while requiring no cold-start fallback.

\section{Condition-Based Analysis}

Breaking the error down by operating condition (Table~\ref{tab:condition}) reveals where the difficulty lies and confirms that all three learned models behave consistently.

\begin{table}[!htbp]
    \centering
    \caption{\label{tab:condition}XGBoost MAE (minutes) by operating condition on route 502.}
    \begin{tabular}{@{}llrr@{}}
        \toprule
        \textbf{Dimension} & \textbf{Condition} & \textbf{N} & \textbf{MAE} \\
        \midrule
        Direction    & Halkapınar $\rightarrow$ Cengizhan & 10{,}957 & 0.380 \\
                     & Cengizhan $\rightarrow$ Halkapınar & 5{,}358  & 0.540 \\
        \midrule
        Time of day  & Morning peak  & 3{,}127 & 0.464 \\
                     & Off-peak      & 9{,}996 & 0.439 \\
                     & Evening peak  & 2{,}387 & 0.388 \\
                     & Night         & 805     & 0.366 \\
        \midrule
        Weather      & Cloudy        & 8{,}978 & 0.418 \\
                     & Clear         & 7{,}013 & 0.444 \\
                     & Rainy         & 324     & 0.595 \\
        \midrule
        Stop zone    & Start (0--33\%)   & 6{,}666 & 0.541 \\
                     & Middle (33--66\%) & 3{,}451 & 0.345 \\
                     & End (66--100\%)   & 6{,}198 & 0.365 \\
        \bottomrule
    \end{tabular}
\end{table}

Three observations stand out. First, the morning peak is the hardest time of day and the night the easiest, as expected from congestion patterns. Second, although weather features were not useful to the models (Chapter~4), the small number of rainy segments that were captured exhibit a clearly higher error ($0.595$ versus about $0.43$ in dry conditions); this is consistent with rain making travel times harder to predict, but the sample is too small ($324$ segments) to train a reliable weather-aware model. Third, and most importantly, the first third of each route (the cold-start zone) has roughly $1.5\times$ the error of the middle and final thirds, confirming that the start of a trip is intrinsically harder because little recent history is available.

\section{Generalization Across Routes}

To verify that the method is not overfitted to a single line, it was applied unchanged to two additional routes, 268 and 565, which have different lengths and traffic profiles. Table~\ref{tab:multiroute} shows that the same pipeline produces consistent, reasonable results on all three routes, with XGBoost MAE ranging from $0.30$ to $0.43$~minutes. This consistency is strong evidence that the approach generalizes rather than exploiting peculiarities of route 502.

\begin{table}[!htbp]
    \centering
    \caption{\label{tab:multiroute}Generalization across three İzmir routes (XGBoost on each route's test set).}
    \begin{tabular}{@{}lrrr@{}}
        \toprule
        \textbf{Route} & \textbf{Segments} & \textbf{MAE (min)} & \textbf{$R^2$} \\
        \midrule
        502 & 81{,}575  & 0.433 & 0.626 \\
        268 & 115{,}803 & 0.371 & 0.542 \\
        565 & 114{,}572 & 0.304 & 0.483 \\
        \bottomrule
    \end{tabular}
\end{table}

\section{Comparison with the Reference Study}

The reference study by Kaya and Kalay \cite{kaya2025} reports an MAE of $2.97$~minutes. The system developed here reaches an MAE of about $0.43$~minutes on route 502, which is numerically much lower. However, this comparison must be made carefully and honestly: the two systems predict at different scales. The reference study predicts arrival times that accumulate over a route, whereas this project predicts the travel time of a single short inter-stop segment (about $1.2$~minutes on average). A smaller absolute error on a smaller quantity does not, by itself, indicate a better model. Consistent with this scale difference, our percentage error (MAPE $\approx 41\%$) and explained variance ($R^2 \approx 0.63$) are less favorable than the reference figures. The honest conclusion is therefore that the project achieves a low absolute error at the segment scale, while the percentage-based metrics reflect the inherent difficulty of predicting very short, highly variable travel times.

\section{Methodological Findings}

Beyond the numerical results, the project produced three methodological findings that are arguably its most valuable outcomes:

\begin{enumerate}
    \item \textbf{A measurement floor governs accuracy.} The achievable MAE of roughly $0.43$~minutes (about $26$~seconds) is close to the precision with which the ground truth can even be measured at a 30-second polling interval. Neither label-precision interpolation, target reframing, nor deep-model tuning moved this floor, which means that better accuracy would require denser positioning data rather than more complex models.
    \item \textbf{Simpler is not worse.} Reducing the feature set, keeping the straightforward target formulation, and using a classical gradient-boosted model all matched or improved upon their more complex alternatives. The strongest single model is the classical XGBoost, not the deep network.
    \item \textbf{Fair evaluation matters.} A rigorous same-test-set comparison overturned an earlier ``deep learning is best'' impression, showing it to be a test-set artifact. This underscores the importance of comparing models on identical data before drawing conclusions.
\end{enumerate}
